%% 09-lifecycle.tex -- generated by tools/build_from_docx.py; edit freely

\hypertarget{lifecycle}{%
\section{Lifecycle}\label{lifecycle}}

An agent that has been admitted is not yet an agent that may act.
Between the acceptance of a descriptor and the first decision that
reaches the network there is a sequence of checks, and this chapter
specifies it. The preceding chapters described contracts an agent
declares; this chapter is where those declarations become verified
bindings, because a declaration the platform never tested is an
assertion by the contributor rather than a property of the system.

The chapter also specifies what happens after activation: the authority
states an agent can occupy, how it moves between them, what the party
that instantiated the agent is told, and what happens when the objective
cannot be met.

\hypertarget{from-descriptor-to-running-agent}{%
\subsection{From Descriptor to Running Agent}\label{from-descriptor-to-running-agent}}

The chain has seven steps, and their order is not arbitrary.

A descriptor is issued for instantiation, together with the values it
leaves to that moment: the target values of objectives bound at
instantiation and the concrete identifiers of its scopes. It goes to the
orchestrator its level designates --- the inter-domain management
function for a global agent, which resolves the service tree into
per-domain portions and issues each to the domain concerned, and the
domain management function directly for a local agent. Resolution
happens once, at instantiation; there is no runtime re-resolution, and a
change is a new agent, for reasons given in Section 9.6. Service
orchestration, in which a customer-facing service is decomposed into
per-domain resource-facing services, runs in parallel and is not part of
this chain.

The descriptor is then admitted, against the checks of Section 9.2. The
domain management function instantiates the tuples and the services
their portion of the service tree declares, assigning identity,
generation, endpoints, credentials and the agent's topic. The activation
gate of Section 9.3 then runs, and only on its outcome is the agent
permitted to emit decisions. The loop then cycles locally, without an
orchestrator round trip per cycle. Assurance is reported throughout to
the requester --- the party that issued the descriptor --- as Section
9.5 describes. At termination, topics are torn down and the generation
retired.

\hypertarget{admission}{%
\subsection{Admission}\label{admission}}

Admission checks the descriptor. Every check below rejects on failure;
none produces a warning that permits the agent through.

The structural checks establish that the descriptor is well-formed and
authentic: a supported schema version, satisfied compatibility, valid
descriptor integrity, artifact attestation where required, and
dependencies available at the required versions.

The semantic checks establish that the declared agent is one the
framework permits. Control scope must be contained within the scope its
observations justify, which must be contained within the declared domain
scope, which must be contained within the trust domain. Every delegated
capability must name a resolvable provider whose declared scope covers
the consumer's, with a pinned interface profile and compatibility range.
Every partial capability must enumerate its subset; an unenumerated
partial declaration is rejected. Insight contracts must contain no
actuation-shaped fields, intent contracts no controller-native command
fields, observation contracts no model outputs. Validity conditions must
be machine-evaluable predicates over observable features. A failure
policy must be present and typed for every implemented role, with no
ambiguous values. The effective loop latency is computed from the
distribution map and the effective latency class stamped. Unverifiable
actions must be declared as such, with authority capped accordingly. A
conflict resolution policy must be present with its priority rules. And
the action space must be bounded: every permitted action carrying
parameter bounds, rate limits and a declared verification method. Table
9-1 states the checks and the condition on which each rejects.

\begin{wptable}{9-1}

\begin{longtable}[]{@{}
  >{\raggedright\arraybackslash}p{(\columnwidth - 4\tabcolsep) * \real{0.1499}}
  >{\raggedright\arraybackslash}p{(\columnwidth - 4\tabcolsep) * \real{0.3499}}
  >{\raggedright\arraybackslash}p{(\columnwidth - 4\tabcolsep) * \real{0.5001}}@{}}
\caption{Admission checks and their rejection conditions}\tabularnewline
\toprule\noalign{}
\begin{minipage}[b]{\linewidth}\raggedright
\textbf{Class}
\end{minipage} & \begin{minipage}[b]{\linewidth}\raggedright
\textbf{Check}
\end{minipage} & \begin{minipage}[b]{\linewidth}\raggedright
\textbf{Rejects when}
\end{minipage} \\
\midrule\noalign{}
\endfirsthead
\toprule\noalign{}
\begin{minipage}[b]{\linewidth}\raggedright
\textbf{Class}
\end{minipage} & \begin{minipage}[b]{\linewidth}\raggedright
\textbf{Check}
\end{minipage} & \begin{minipage}[b]{\linewidth}\raggedright
\textbf{Rejects when}
\end{minipage} \\
\midrule\noalign{}
\endhead
\bottomrule\noalign{}
\endlastfoot
\textbf{Structural} & schema\_\allowbreak{}version supported;
compatibility satisfied & Unsupported or incompatible version. \\
\textbf{Structural} & Descriptor integrity and artifact attestation &
Invalid hash or signature; attestation absent where required. \\
\textbf{Structural} & Dependencies resolvable at required versions &
Dependency absent or out of range. \\
\textbf{Semantic} & Scope containment & control\_\allowbreak{}scope not
within observation-justified scope, declared domain scope or trust
domain. \\
\textbf{Semantic} & Delegated capability resolution & Provider
unresolvable, scope not covering the consumer, or interface profile
unpinned. \\
\textbf{Semantic} & Partial capability enumeration & Subset not
enumerated. \\
\textbf{Semantic} & Contract purity & Actuation-shaped fields in
insight; controller-native commands in intent; model outputs in
observation. \\
\textbf{Semantic} & Validity conditions machine-evaluable & Condition
not a predicate over observable features. \\
\textbf{Semantic} & Failure policy present and typed & Absent for an
implemented role, or ambiguous value. \\
\textbf{Semantic} & Effective latency class computable & Distribution
map insufficient to compute the class. \\
\textbf{Semantic} & Unverifiable actions declared & Undeclared
unverifiable action, or authority above the permitted ceiling. \\
\textbf{Semantic} & Conflict resolution policy present & Absent, or
priority rules missing. \\
\textbf{Semantic} & Action space bounded & Any permitted action lacking
parameter bounds, rate limits or a verification method. \\
\end{longtable}

\end{wptable}

The consistent shape of these checks is that the platform is flexible
about what may be declared and strict about requiring that it be
declared. A contributor may describe an unusual agent; a contributor may
not describe an agent incompletely and have the gaps filled by platform
defaults, because a default the contributor never saw is a behavior
nobody has taken responsibility for.

\hypertarget{the-activation-gate}{%
\subsection{The Activation Gate}\label{the-activation-gate}}

A running container proves nothing. Placement establishes that an agent
exists; it establishes nothing about whether the agent can observe what
it claims to observe, reach what it claims to control, or have its
actions verified. The activation gate is where each of those is proven,
and decision emission is enabled last.

Nine checks run in order, operated by the domain management function
against platform-verified evidence rather than agent self-report. Table
9-2 states them in order with the evidence each requires.

Live observation: the Monitoring System is publishing within scope, with
coverage at or above the declared minimum. Not that the Monitoring
System started, but that data has been observed. Model availability and
compatibility: declared artifacts and endpoints resolvable and within
pinned version ranges. Actuation authority: control scope and action
space registered and consistent with the authority granted at admission
within what the descriptor declares. Controller reachability: the
southbound controller reachable and its interface version inside the
pinned range. Verification path: a read-back route proven for every
action declared verifiable, with actions failing this either demoted to
unverifiable and authority-capped, or rejected. Rollback capability: a
rollback route proven wherever the failure policy declares one. Conflict
binding: the mediation endpoint bound and responding, which is what
makes the prohibition on self-approval non-bypassable rather than merely
stated. Feedback route: the actuation feedback path proven to both the
Decision Engine and the domain management function.

\begin{wptable}{9-2}

\begin{longtable}[]{@{}
  >{\raggedright\arraybackslash}p{(\columnwidth - 4\tabcolsep) * \real{0.0625}}
  >{\raggedright\arraybackslash}p{(\columnwidth - 4\tabcolsep) * \real{0.2625}}
  >{\raggedright\arraybackslash}p{(\columnwidth - 4\tabcolsep) * \real{0.6750}}@{}}
\caption{The nine activation gate checks, in order}\tabularnewline
\toprule\noalign{}
\begin{minipage}[b]{\linewidth}\raggedright
\textbf{\#}
\end{minipage} & \begin{minipage}[b]{\linewidth}\raggedright
\textbf{Check}
\end{minipage} & \begin{minipage}[b]{\linewidth}\raggedright
\textbf{Evidence required}
\end{minipage} \\
\midrule\noalign{}
\endfirsthead
\toprule\noalign{}
\begin{minipage}[b]{\linewidth}\raggedright
\textbf{\#}
\end{minipage} & \begin{minipage}[b]{\linewidth}\raggedright
\textbf{Check}
\end{minipage} & \begin{minipage}[b]{\linewidth}\raggedright
\textbf{Evidence required}
\end{minipage} \\
\midrule\noalign{}
\endhead
\bottomrule\noalign{}
\endlastfoot
\textbf{1} & Live observation & MS publishing within scope,
scope\_\allowbreak{}coverage at or above the declared minimum. Data
observed, not process started. \\
\textbf{2} & Model availability and compatibility & Declared artifacts
and endpoints resolvable, within pinned version ranges. \\
\textbf{3} & Actuation authority & Control scope and action space
registered and consistent with the RFS-granted authority. \\
\textbf{4} & Controller reachability & Southbound controller reachable,
interface version inside the pinned range. \\
\textbf{5} & Verification path & Read-back route proven for every action
declared verifiable. Failure demotes the action to unverifiable with
authority capped, or rejects it. \\
\textbf{6} & Rollback capability & Rollback route proven wherever the
failure policy declares one. \\
\textbf{7} & Conflict binding & Mediation endpoint bound and
responding. \\
\textbf{8} & Feedback route & Actuation feedback path proven to both the
DE and the DMO. \\
\textbf{9} & Correlated dry-run & One end-to-end cycle completed with
ACT inhibited, feedback correlated. \\
\end{longtable}

\end{wptable}

Every one of these is a precondition for a safe decision. Enabling
decisions before they hold means the first decision of a new agent's
life is taken under unverified conditions, which is the single
highest-risk moment in the agent lifecycle. Gate failure leaves the
agent in admitted\_\allowbreak{}not\_\allowbreak{}active rather than
partially enabled.

\hypertarget{the-correlated-dry-run-cycle}{%
\subsubsection{The correlated dry-run cycle}\label{the-correlated-dry-run-cycle}}

A binding that carries executable intent is tested with actuation
inhibited before it is activated, and acceptance requires one correlated
end-to-end dry-run cycle. This is the platform-level proof that the
whole chain is operational: observation through insight, decision,
inhibited actuation and feedback correlation, exercised once with the
network untouched.

The dry-run is what distinguishes eight checks that each pass
individually from a chain that actually carries a cycle end to end.
Components can be severally reachable and still fail to correlate.

\hypertarget{gate-outcomes}{%
\subsubsection{Gate outcomes}\label{gate-outcomes}}

The gate is not binary. It permits full, restricted, advisory or
observation-only operation, and these map onto the authority states of
the next section: full to closed\_\allowbreak{}loop\_\allowbreak{}full,
restricted to closed\_\allowbreak{}loop\_\allowbreak{}reduced, advisory
to advisory, in which intents are published as recommendations and
proceed only when approved, and observation-only to
monitoring\_\allowbreak{}only.\allowbreak{}

Advisory is the correct outcome for a newly admitted experimental agent,
and is the adopted entry state for the extension path. An agent that
must prove itself is better placed where its decisions are visible and
inert than excluded until someone is confident enough to grant it
authority, because the second arrangement produces no evidence on which
confidence could be based.

\hypertarget{authority-states}{%
\subsection{Authority States}\label{authority-states}}

An agent occupies exactly one of the seven authority states of Table
3-12, held in the registry. In
admitted\_\allowbreak{}not\_\allowbreak{}active, admission has passed
and the gate is unsatisfied, and no decision is emitted. In
monitoring\_\allowbreak{}only, the Monitoring System and Analytics
Engine are active and the Decision Engine is not invoked. In advisory,
intents are published as recommendations, and one proceeds only when
approved, as an intent carrying the approval reference and after
mediation; this is the entry state of experimental agents. In
closed\_\allowbreak{}loop\_\allowbreak{}reduced, decisions are enabled
within a narrowed control scope, tighter bounds and a lower priority
band, which is also the state of agents downgraded by evaluation. In
closed\_\allowbreak{}loop\_\allowbreak{}full, the agent acts within its
full granted authority. In suspended, decision emission is withdrawn,
actuation is inhibited except for completing verification and permitted
rollback of transactions in flight, and monitoring and analytics may
continue. The state terminated is terminal. Table 9-3 states the states
and the transitions permitted between them.

\begin{wptable}{9-3}

\begin{longtable}[]{@{}
  >{\raggedright\arraybackslash}p{(\columnwidth - 4\tabcolsep) * \real{0.2750}}
  >{\raggedright\arraybackslash}p{(\columnwidth - 4\tabcolsep) * \real{0.3875}}
  >{\raggedright\arraybackslash}p{(\columnwidth - 4\tabcolsep) * \real{0.3375}}@{}}
\caption{Authority states and their permitted
transitions}\tabularnewline
\toprule\noalign{}
\begin{minipage}[b]{\linewidth}\raggedright
\textbf{State}
\end{minipage} & \begin{minipage}[b]{\linewidth}\raggedright
\textbf{Decision emission and actuation}
\end{minipage} & \begin{minipage}[b]{\linewidth}\raggedright
\textbf{Transitions}
\end{minipage} \\
\midrule\noalign{}
\endfirsthead
\toprule\noalign{}
\begin{minipage}[b]{\linewidth}\raggedright
\textbf{State}
\end{minipage} & \begin{minipage}[b]{\linewidth}\raggedright
\textbf{Decision emission and actuation}
\end{minipage} & \begin{minipage}[b]{\linewidth}\raggedright
\textbf{Transitions}
\end{minipage} \\
\midrule\noalign{}
\endhead
\bottomrule\noalign{}
\endlastfoot
\textbf{admitted\_\allowbreak{}not\_\allowbreak{}active} & No decision
emitted; actuation inhibited, dry run only during the activation gate. &
To the state of the gate outcome; to terminated. \\
\textbf{monitoring\_\allowbreak{}only} & DE not invoked; actuation
inhibited. & To a higher state within the ceiling on gate outcome or
promotion; to suspended; to terminated. \\
\textbf{advisory} & Intents published as recommendations; an approved
recommendation proceeds as an intent carrying the approval reference,
after mediation. & To closed\_\allowbreak{}loop\_\allowbreak{}reduced on
sustained conformant evaluation; to a lower state automatically on
declared thresholds; to suspended. \\
\textbf{closed\_\allowbreak{}loop\_\allowbreak{}reduced} & Intents
within a narrowed control scope, tightened bounds and a lower priority
band. & To closed\_\allowbreak{}loop\_\allowbreak{}full on sustained
conformance; to a lower state automatically on declared thresholds; to
suspended. \\
\textbf{closed\_\allowbreak{}loop\_\allowbreak{}full} & Intents within
the full granted authority. & To a lower state automatically on declared
thresholds; to suspended. \\
\textbf{suspended} & Decision emission withdrawn; actuation inhibited
except completion of verification and permitted rollback of transactions
in flight; MS and AE may continue. & To a state within the ceiling on a
decided act with a record; to terminated. \\
\textbf{terminated} & None. & Terminal. \\
\end{longtable}

\end{wptable}

Promotion requires the gate, and for experimental agents sustained
conformant evaluation as well. Demotion is automatic on declared
thresholds and requires no human step. Every transition is a lifecycle
event to the domain management function and appears in assurance.

The states are deliberately reversible and deliberately explicit. An
authority that can only be granted, never withdrawn without human
intervention, is an authority that persists through exactly the
conditions that should remove it.

\hypertarget{per-instance-isolation}{%
\subsubsection{Per-instance isolation}\label{per-instance-isolation}}

The unit of isolation is the agent instance, identified by the triple of
agent, assigned domain and instance. Each instance holds separate policy
state, model state, decision history, action correlation, evaluation
windows and lifecycle. Shared physical resources do not imply a shared
agent. A read-only Monitoring System collector may be shared only where
the service trees of the instances it serves declare it shared and
read-only, with separately scoped logical streams per instance.

Generation fencing is per agent instance, and the identifier triple is
what prevents state collision when several agent instances occupy one
domain simultaneously.

The capabilities realized in a domain are those the descriptors
instantiated there declare, so a closed-loop agent in one domain may
coexist with a monitoring agent in another. The monitoring flag of a
resource-facing service provisions a monitoring capability of the
service orchestration pipeline; it does not imply that any agent exists,
and a closed loop exists only where a descriptor declaring one is
instantiated.

\hypertarget{assurance-to-the-requester}{%
\subsection{Assurance to the Requester}\label{assurance-to-the-requester}}

The requester --- the party that issued the descriptor for instantiation
--- must be able to learn that the agent is meeting its objectives and
fully handling what it acts on, without inspecting agent internals. Five
things are reported.

Objective status against the agent's declared objectives, as met,
at-risk or not-met. An actuation outcome summary, which includes counts
of acknowledged\_\allowbreak{}unverified and unresolved outcomes.
Authority-state changes, meaning any downgrade or suspension of a
participating agent. Unachievable-objective events, raised when a
Decision Engine exhausts its re-decision budget. And unresolved-conflict
and deferral rates.

The second of these carries the weight. Acknowledged-unverified and
unresolved are precisely the outcomes that would otherwise present as
success, and an agent that reports healthy while it achieves nothing
physically is the exact failure this reporting exists to expose.
Assurance is assembled by the domain management function from
platform-emitted lifecycle and feedback events, not from agent
self-reports, for the same reason the activation gate is: a component
reporting on its own health reports the health it believes it has.

The assurance vocabulary is adopted rather than invented. States are
healthy, observe, degraded, restricted, unsafe and failed;
recommendations are continue, tune, scale, restart-component,
suspend-actions and terminate. The evaluator provides evidence and
urgency and does not execute infrastructure or network changes; the
owning domain or inter-domain management function applies the authorized
lifecycle operation. This is the same separation Chapter 5 reached from
the analytics side, and the two agree. Table 9-4 states the assurance
states and the recommendations associated with each.

\begin{wptable}{9-4}

\begin{longtable}[]{@{}
  >{\raggedright\arraybackslash}p{(\columnwidth - 4\tabcolsep) * \real{0.1625}}
  >{\raggedright\arraybackslash}p{(\columnwidth - 4\tabcolsep) * \real{0.5000}}
  >{\raggedright\arraybackslash}p{(\columnwidth - 4\tabcolsep) * \real{0.3375}}@{}}
\caption{Assurance states and their associated
recommendations}\tabularnewline
\toprule\noalign{}
\begin{minipage}[b]{\linewidth}\raggedright
\textbf{Assurance state}
\end{minipage} & \begin{minipage}[b]{\linewidth}\raggedright
\textbf{Meaning}
\end{minipage} & \begin{minipage}[b]{\linewidth}\raggedright
\textbf{Recommendations}
\end{minipage} \\
\midrule\noalign{}
\endfirsthead
\toprule\noalign{}
\begin{minipage}[b]{\linewidth}\raggedright
\textbf{Assurance state}
\end{minipage} & \begin{minipage}[b]{\linewidth}\raggedright
\textbf{Meaning}
\end{minipage} & \begin{minipage}[b]{\linewidth}\raggedright
\textbf{Recommendations}
\end{minipage} \\
\midrule\noalign{}
\endhead
\bottomrule\noalign{}
\endlastfoot
\textbf{HEALTHY} & Objective met, no adverse indicators. & CONTINUE \\
\textbf{OBSERVE} & Indicators warrant closer watching. & CONTINUE,
TUNE \\
\textbf{DEGRADED} & Objective at risk or partially unmet. & TUNE,
SCALE \\
\textbf{RESTRICTED} & Operating below declared authority. & TUNE,
RESTART\_\allowbreak{}COMPONENT \\
\textbf{UNSAFE} & Conditions under which continued action is not
warranted. & SUSPEND\_\allowbreak{}ACTIONS \\
\textbf{FAILED} & Objective not met, no path to recovery within the
binding. & TERMINATE \\
\end{longtable}

\end{wptable}

\hypertarget{replacement-runs-in-parallel}{%
\subsection{Replacement Runs in Parallel}\label{replacement-runs-in-parallel}}

Where an agent cannot meet its objective, it reports degraded, and a
change of objective, scope or service tree is made by instantiating a
new agent from a new descriptor; for a global agent, the same holds
where its per-domain portions cannot meet the end-to-end objective.
There is no fast runtime re-planner, because a fast global re-planner
would be a global Decision Engine under another name, holding
cross-domain decision authority that the federation model places with
the inter-domain management function and nowhere else.

The absence of a runtime re-planner does not mean the objective is
unserved while the replacement is prepared. The incumbent continues
operating throughout: it keeps observing, keeps deciding within its
granted authority and keeps actuating, at the degraded objective it is
able to meet. The replacement is admitted, instantiated and gated in
parallel, against the same checks as any other. What the lifecycle path
costs is therefore the cutover, not the whole preparation.

\hypertarget{handover}{%
\subsubsection{Handover}\label{handover}}

An incumbent and its replacement must never hold authority over the same
action space at the same time. The incumbent instances retain authority
until the replacement instances have satisfied the activation gate; only
then is authority withdrawn from the incumbents and the replacements
activated. Generation fencing on the agent instance is the mechanism:
intents carrying a retired generation are refused at the conflict
fabric, so an incumbent that has not yet observed its own demotion
cannot act after the handover point.

Where a replacement instance fails its gate, the incumbents are not
demoted and continue, degraded. A handover that demoted the incumbents
first would convert a failed replacement into an unserved objective.

Should the consortium later require adaptation faster than a gated
handover allows, its natural home is a supporting orchestration
application operating between domains, as the architecture work
envisages. That is a future decision and is deliberately not designed
into the framework now.

\hypertarget{termination}{%
\subsection{Termination}\label{termination}}

Termination retires the generation, tears down the agent's topics and
their subscriptions, and releases the identity. An agent in termination
completes or abandons its outstanding transactions according to its
declared failure behavior; it does not acquire a final grace period in
which its actions are unmediated.

Outstanding actuation transactions unresolved at termination are
reported as unresolved, not silently closed. A terminated agent that
leaves the network in a state nobody recorded is worse than one that
leaves a known fault behind it.

\hypertarget{requirements-6}{%
\subsection{Requirements}\label{requirements-6}}

\begin{requirement}{Req-T23-LC-01}{MUST}{Admission rejects on failure}
\reqfield{Description / rationale}{%
Every admission check, structural and semantic, shall reject the
descriptor on failure. No admission check shall produce a warning that
permits the agent to proceed, and no unstated element of a descriptor
shall be supplied by a platform default.}
\reqfield{Dependency}{%
Req-T23-FW-06}
\reqfield{Traceability (Forward)}{%
\begin{itemize}
\tightlist
\item
  Checked at admission by the IDMO or DMO.
\end{itemize}}
\reqfield{UNITY-6G components}{%
\begin{itemize}
\tightlist
\item
  Admission (IDMO/DMO)
\end{itemize}}
\reqfield{Linked KPIs}{%
\begin{itemize}
\tightlist
\item
  KPI\#36: Guarantee 99.99999\% service availability and reliability
  based on effective monitoring and LCM.
\item
  KPI\#1: Deliver four distributed management entities (MS, AE, DE and
  ACT) deployable in each heterogeneous domain.
\end{itemize}}
\end{requirement}

\begin{requirement}{Req-T23-LC-02}{MUST}{Decision emission is enabled last}
\reqfield{Description / rationale}{%
The platform shall not enable decision emission for an agent instance
until every activation gate check has been satisfied. Gate failure shall
leave the instance in admitted\_\allowbreak{}not\_\allowbreak{}active,
never in a partially enabled state.}
\reqfield{Dependency}{%
Req-T23-LC-01}
\reqfield{Traceability (Forward)}{%
\begin{itemize}
\tightlist
\item
  Gate operated by the DMO before emission is enabled.
\end{itemize}}
\reqfield{UNITY-6G components}{%
\begin{itemize}
\tightlist
\item
  DMO
\item
  Platform runtime
\end{itemize}}
\reqfield{Linked KPIs}{%
\begin{itemize}
\tightlist
\item
  KPI\#1: Deliver four distributed management entities (MS, AE, DE and
  ACT) deployable in each heterogeneous domain.
\item
  KPI\#21: Reduce the performance penalty due to conflicts by 20\% by
  means of DLT, hierarchical algorithms and policy-based conflict
  resolution schemes.
\end{itemize}}
\end{requirement}

\begin{requirement}{Req-T23-LC-03}{MUST}{The gate runs on platform-verified evidence}
\reqfield{Description / rationale}{%
Every gate check shall be satisfied against evidence the platform
observes, and never against agent self-report. Live observation shall
mean data observed within scope at or above the declared coverage
minimum, not that the MS process has started.}
\reqfield{Dependency}{%
Req-T23-LC-02}
\reqfield{Traceability (Forward)}{%
\begin{itemize}
\tightlist
\item
  Evidence collected by the platform runtime.
\end{itemize}}
\reqfield{UNITY-6G components}{%
\begin{itemize}
\tightlist
\item
  DMO
\item
  Platform runtime
\end{itemize}}
\reqfield{Linked KPIs}{%
\begin{itemize}
\tightlist
\item
  KPI\#1: Deliver four distributed management entities (MS, AE, DE and
  ACT) deployable in each heterogeneous domain.
\item
  KPI\#36: Guarantee 99.99999\% service availability and reliability
  based on effective monitoring and LCM.
\end{itemize}}
\end{requirement}

\begin{requirement}{Req-T23-LC-04}{MUST}{A correlated dry-run cycle precedes activation}
\reqfield{Description / rationale}{%
A binding carrying executable intent shall complete one correlated
end-to-end cycle with actuation inhibited before activation, exercising
observation, insight, decision, inhibited actuation and feedback
correlation. Newly admitted experimental agents shall enter at advisory.}
\reqfield{Dependency}{%
Req-T23-LC-02}
\reqfield{Traceability (Forward)}{%
\begin{itemize}
\tightlist
\item
  Dry-run cycle correlated by the platform before activation.
\end{itemize}}
\reqfield{UNITY-6G components}{%
\begin{itemize}
\tightlist
\item
  DMO
\item
  Platform runtime
\end{itemize}}
\reqfield{Linked KPIs}{%
\begin{itemize}
\tightlist
\item
  KPI\#41: Reduction of 70\% of deployment time of AI-based
  resource/energy management algorithms.
\item
  KPI\#1: Deliver four distributed management entities (MS, AE, DE and
  ACT) deployable in each heterogeneous domain.
\end{itemize}}
\end{requirement}

\begin{requirement}{Req-T23-LC-05}{MUST}{Authority state is registry-held and singular}
\reqfield{Description / rationale}{%
An agent instance shall occupy exactly one authority state, held in the
registry. The intent-egress gate shall consult the registry authority
state per cycle, and authority shall never be inferred from the
descriptor alone.}
\reqfield{Dependency}{%
Req-T23-FW-23}
\reqfield{Traceability (Forward)}{%
\begin{itemize}
\tightlist
\item
  Consulted per cycle at intent egress.
\end{itemize}}
\reqfield{UNITY-6G components}{%
\begin{itemize}
\tightlist
\item
  DMO
\item
  Conflict fabric
\end{itemize}}
\reqfield{Linked KPIs}{%
\begin{itemize}
\tightlist
\item
  KPI\#21: Reduce the performance penalty due to conflicts by 20\% by
  means of DLT, hierarchical algorithms and policy-based conflict
  resolution schemes.
\item
  KPI\#1: Deliver four distributed management entities (MS, AE, DE and
  ACT) deployable in each heterogeneous domain.
\end{itemize}}
\end{requirement}

\begin{requirement}{Req-T23-LC-06}{MUST}{Demotion is automatic, promotion is earned}
\reqfield{Description / rationale}{%
Demotion shall occur automatically on declared thresholds and shall
require no human step. Promotion shall require gate satisfaction and,
for experimental agents, sustained conformant evaluation. Every
transition shall be emitted as a lifecycle event and shall appear in
assurance.}
\reqfield{Dependency}{%
Req-T23-LC-05}
\reqfield{Traceability (Forward)}{%
\begin{itemize}
\tightlist
\item
  Thresholds evaluated by the evaluation path.
\item
  Transitions emitted as lifecycle events.
\end{itemize}}
\reqfield{UNITY-6G components}{%
\begin{itemize}
\tightlist
\item
  DMO
\item
  Evaluation
\end{itemize}}
\reqfield{Linked KPIs}{%
\begin{itemize}
\tightlist
\item
  KPI\#21: Reduce the performance penalty due to conflicts by 20\% by
  means of DLT, hierarchical algorithms and policy-based conflict
  resolution schemes.
\item
  KPI\#31: Reduce service recovery time below 180 s.
\end{itemize}}
\end{requirement}

\begin{requirement}{Req-T23-LC-07}{MUST}{Isolation is per intent, domain and instance}
\reqfield{Description / rationale}{%
The unit of isolation shall be the triple of intent, assigned domain and
agent instance, each holding separate policy state, model state,
decision history, action correlation, evaluation windows and lifecycle.
Generation fencing shall be per agent instance. A shared read-only MS
collector shall be permitted only where the RFS declares it explicitly,
with separately scoped logical streams per intent.}
\reqfield{Dependency}{%
Req-T23-FW-12}
\reqfield{Traceability (Forward)}{%
\begin{itemize}
\tightlist
\item
  Identifier triple assigned at placement.
\end{itemize}}
\reqfield{UNITY-6G components}{%
\begin{itemize}
\tightlist
\item
  IDMO
\item
  DMO
\end{itemize}}
\reqfield{Linked KPIs}{%
\begin{itemize}
\tightlist
\item
  KPI\#26: Increase the ratio of service management and LCM tasks
  resolved by local AE/DE components.
\item
  KPI\#1: Deliver four distributed management entities (MS, AE, DE and
  ACT) deployable in each heterogeneous domain.
\end{itemize}}
\end{requirement}

\begin{requirement}{Req-T23-LC-08}{MUST}{Assurance reaches the requester}
\reqfield{Description / rationale}{%
The platform shall report to the requester, without exposing agent
internals: objective status against the agent\textquotesingle s declared
objectives; an actuation outcome summary including counts of
acknowledged\_\allowbreak{}unverified and unresolved; authority-state
changes; unachievable-objective events; and unresolved-conflict and
deferral rates.}
\reqfield{Dependency}{%
Req-T23-ACT-09}
\reqfield{Traceability (Forward)}{%
\begin{itemize}
\tightlist
\item
  Assembled by the DMO, aggregated E2E by the IDMO.
\end{itemize}}
\reqfield{UNITY-6G components}{%
\begin{itemize}
\tightlist
\item
  DMO
\item
  IDMO
\item
  Intent layer
\end{itemize}}
\reqfield{Linked KPIs}{%
\begin{itemize}
\tightlist
\item
  KPI\#36: Guarantee 99.99999\% service availability and reliability
  based on effective monitoring and LCM.
\item
  KPI\#21: Reduce the performance penalty due to conflicts by 20\% by
  means of DLT, hierarchical algorithms and policy-based conflict
  resolution schemes.
\end{itemize}}
\end{requirement}

\begin{requirement}{Req-T23-LC-09}{MUST}{Evaluation recommends, the orchestrator enacts}
\reqfield{Description / rationale}{%
The evaluation path shall provide evidence, an assurance state and a
recommendation, and shall not execute infrastructure or network changes.
The owning DMO or IDMO shall apply the authorized lifecycle operation.}
\reqfield{Dependency}{%
Req-T23-AE-06}
\reqfield{Traceability (Forward)}{%
\begin{itemize}
\tightlist
\item
  Separation enforced at the lifecycle operation interface.
\end{itemize}}
\reqfield{UNITY-6G components}{%
\begin{itemize}
\tightlist
\item
  Evaluation
\item
  DMO
\item
  IDMO
\end{itemize}}
\reqfield{Linked KPIs}{%
\begin{itemize}
\tightlist
\item
  KPI\#21: Reduce the performance penalty due to conflicts by 20\% by
  means of DLT, hierarchical algorithms and policy-based conflict
  resolution schemes.
\item
  KPI\#36: Guarantee 99.99999\% service availability and reliability
  based on effective monitoring and LCM.
\end{itemize}}
\end{requirement}

\begin{requirement}{Req-T23-LC-10}{MUST}{Termination leaves no silent state}
\reqfield{Description / rationale}{%
Termination shall retire the generation, tear down the
agent\textquotesingle s topics and their subscriptions and release the
identity. Outstanding actuation transactions unresolved at termination
shall be reported as unresolved and shall not be silently closed. A
terminating agent shall acquire no grace period in which its actions are
unmediated.}
\reqfield{Dependency}{%
Req-T23-ACT-10}
\reqfield{Traceability (Forward)}{%
\begin{itemize}
\tightlist
\item
  Teardown driven by the service lifecycle.
\item
  Unresolved transactions reported to assurance.
\end{itemize}}
\reqfield{UNITY-6G components}{%
\begin{itemize}
\tightlist
\item
  DMO
\item
  SMB
\item
  Platform runtime
\end{itemize}}
\reqfield{Linked KPIs}{%
\begin{itemize}
\tightlist
\item
  KPI\#31: Reduce service recovery time below 180 s.
\item
  KPI\#36: Guarantee 99.99999\% service availability and reliability
  based on effective monitoring and LCM.
\end{itemize}}
\end{requirement}

\begin{requirement}{Req-T23-LC-11}{MUST}{Replacement runs in parallel with the incumbent}
\reqfield{Description / rationale}{%
Where an agent cannot meet its objective, it shall report degraded, and
a change of objective, scope or service tree shall require a new agent
instantiated from a new descriptor. The platform shall provide no fast
runtime re-planner holding cross-domain decision authority. Incumbent
instances shall continue operating while the replacement is admitted,
instantiated and gated in parallel.}
\reqfield{Dependency}{%
Req-T23-DE-10}
\reqfield{Traceability (Forward)}{%
\begin{itemize}
\tightlist
\item
  Enforced by the absence of a runtime re-resolution path.
\end{itemize}}
\reqfield{UNITY-6G components}{%
\begin{itemize}
\tightlist
\item
  IDMO
\item
  DMO
\end{itemize}}
\reqfield{Linked KPIs}{%
\begin{itemize}
\tightlist
\item
  KPI\#21: Reduce the performance penalty due to conflicts by 20\% by
  means of DLT, hierarchical algorithms and policy-based conflict
  resolution schemes.
\item
  KPI\#26: Increase the ratio of service management and LCM tasks
  resolved by local AE/DE components.
\end{itemize}}
\end{requirement}

\begin{requirement}{Req-T23-LC-12}{MUST}{Handover is gated, and never overlaps authority}
\reqfield{Description / rationale}{%
Two decompositions shall never hold authority over the same action space
simultaneously. Incumbent instances shall retain authority until the
replacement instances have satisfied the activation gate, at which point
authority shall be withdrawn from the incumbents and the replacements
activated. Intents carrying a retired generation shall be refused at the
conflict fabric. Where a replacement instance fails its gate, the
incumbents shall not be demoted.}
\reqfield{Dependency}{%
Req-T23-LC-11}
\reqfield{Traceability (Forward)}{%
\begin{itemize}
\tightlist
\item
  Generation fencing enforced at the conflict fabric.
\item
  Gate satisfaction precedes incumbent demotion.
\end{itemize}}
\reqfield{UNITY-6G components}{%
\begin{itemize}
\tightlist
\item
  IDMO
\item
  DMO
\item
  Conflict fabric
\end{itemize}}
\reqfield{Linked KPIs}{%
\begin{itemize}
\tightlist
\item
  KPI\#31: Reduce service recovery time below 180 s.
\item
  KPI\#21: Reduce the performance penalty due to conflicts by 20\% by
  means of DLT, hierarchical algorithms and policy-based conflict
  resolution schemes.
\end{itemize}}
\end{requirement}

\hypertarget{contribution-to-the-project-kpis-5}{%
\subsection{CONTRIBUTION TO THE PROJECT KPIs}\label{contribution-to-the-project-kpis-5}}

Table 9-5 states what the lifecycle contributes to each KPI it bears on,
and Table 9-6 maps the requirements of this chapter onto those KPIs.

\begin{wptable}{9-5}

\begin{longtable}[]{@{}
  >{\raggedright\arraybackslash}p{(\columnwidth - 4\tabcolsep) * \real{0.1499}}
  >{\raggedright\arraybackslash}p{(\columnwidth - 4\tabcolsep) * \real{0.6375}}
  >{\raggedright\arraybackslash}p{(\columnwidth - 4\tabcolsep) * \real{0.2126}}@{}}
\caption{Contribution of the lifecycle to the project
KPIs}\tabularnewline
\toprule\noalign{}
\begin{minipage}[b]{\linewidth}\raggedright
\textbf{KPI}
\end{minipage} & \begin{minipage}[b]{\linewidth}\raggedright
\textbf{Contribution of the lifecycle}
\end{minipage} & \begin{minipage}[b]{\linewidth}\raggedright
\textbf{Requirements}
\end{minipage} \\
\midrule\noalign{}
\endfirsthead
\toprule\noalign{}
\begin{minipage}[b]{\linewidth}\raggedright
\textbf{KPI}
\end{minipage} & \begin{minipage}[b]{\linewidth}\raggedright
\textbf{Contribution of the lifecycle}
\end{minipage} & \begin{minipage}[b]{\linewidth}\raggedright
\textbf{Requirements}
\end{minipage} \\
\midrule\noalign{}
\endhead
\bottomrule\noalign{}
\endlastfoot
\textbf{KPI\#1} & The gate turns declared contracts into verified
bindings, which is what lets heterogeneous contributed agents share one
platform. & Req-T23-LC-02, Req-T23-LC-03 \\
\textbf{KPI\#21} & The conflict binding is proven before decision
emission is enabled, the registry authority state is consulted at every
intent egress, and a gated handover never lets two decompositions hold
authority over the same action space at once. & Req-T23-LC-02,
Req-T23-LC-05, Req-T23-LC-12 \\
\textbf{KPI\#26} & Per-instance isolation on the agent, domain and
instance triple prevents state collision where several agent instances
occupy one domain. & Req-T23-LC-07 \\
\textbf{KPI\#31} & Automatic demotion withdraws authority without a
human step; gated handover keeps the incumbent decomposition serving
while a replacement is prepared, so recovery costs the cutover rather
than the whole lifecycle path. & Req-T23-LC-06, Req-T23-LC-10,
Req-T23-LC-12 \\
\textbf{KPI\#36} & Admission rejects on failure rather than warning, so
an incompletely declared agent never reaches placement, and assurance is
assembled from platform-emitted events and surfaces the outcomes that
would otherwise present as success. & Req-T23-LC-01, Req-T23-LC-03,
Req-T23-LC-08, Req-T23-LC-09, Req-T23-LC-10 \\
\textbf{KPI\#41} & Advisory entry lets a new AI-based contribution be
deployed at once and produce evidence, rather than waiting outside the
platform until it has earned authority. & Req-T23-LC-04 \\
\end{longtable}

\end{wptable}

\begin{wptable}{9-6}\footnotesize

\begin{longtable}[]{@{}
  >{\raggedright\arraybackslash}p{(\columnwidth - 12\tabcolsep) * \real{0.1813}}
  >{\raggedright\arraybackslash}p{(\columnwidth - 12\tabcolsep) * \real{0.1365}}
  >{\raggedright\arraybackslash}p{(\columnwidth - 12\tabcolsep) * \real{0.1365}}
  >{\raggedright\arraybackslash}p{(\columnwidth - 12\tabcolsep) * \real{0.1365}}
  >{\raggedright\arraybackslash}p{(\columnwidth - 12\tabcolsep) * \real{0.1365}}
  >{\raggedright\arraybackslash}p{(\columnwidth - 12\tabcolsep) * \real{0.1365}}
  >{\raggedright\arraybackslash}p{(\columnwidth - 12\tabcolsep) * \real{0.1363}}@{}}
\caption{Mapping of lifecycle requirements onto the project
KPIs}\tabularnewline
\toprule\noalign{}
\begin{minipage}[b]{\linewidth}\raggedright
\textbf{KPI\textbackslash{}}

\textbf{Requirement}
\end{minipage} & \begin{minipage}[b]{\linewidth}\raggedright
\textbf{\#1}
\end{minipage} & \begin{minipage}[b]{\linewidth}\raggedright
\textbf{\#21}
\end{minipage} & \begin{minipage}[b]{\linewidth}\raggedright
\textbf{\#26}
\end{minipage} & \begin{minipage}[b]{\linewidth}\raggedright
\textbf{\#31}
\end{minipage} & \begin{minipage}[b]{\linewidth}\raggedright
\textbf{\#36}
\end{minipage} & \begin{minipage}[b]{\linewidth}\raggedright
\textbf{\#41}
\end{minipage} \\
\midrule\noalign{}
\endfirsthead
\toprule\noalign{}
\begin{minipage}[b]{\linewidth}\raggedright
\textbf{KPI\textbackslash{}}

\textbf{Requirement}
\end{minipage} & \begin{minipage}[b]{\linewidth}\raggedright
\textbf{\#1}
\end{minipage} & \begin{minipage}[b]{\linewidth}\raggedright
\textbf{\#21}
\end{minipage} & \begin{minipage}[b]{\linewidth}\raggedright
\textbf{\#26}
\end{minipage} & \begin{minipage}[b]{\linewidth}\raggedright
\textbf{\#31}
\end{minipage} & \begin{minipage}[b]{\linewidth}\raggedright
\textbf{\#36}
\end{minipage} & \begin{minipage}[b]{\linewidth}\raggedright
\textbf{\#41}
\end{minipage} \\
\midrule\noalign{}
\endhead
\bottomrule\noalign{}
\endlastfoot
\textbf{Req-T23-LC-01} & \kpimark{} & & & & \kpimark{} & \\
\textbf{Req-T23-LC-02} & \kpimark{} & \kpimark{} & & & & \\
\textbf{Req-T23-LC-03} & \kpimark{} & & & & \kpimark{} & \\
\textbf{Req-T23-LC-04} & \kpimark{} & & & & & \kpimark{} \\
\textbf{Req-T23-LC-05} & \kpimark{} & \kpimark{} & & & & \\
\textbf{Req-T23-LC-06} & & \kpimark{} & & \kpimark{} & & \\
\textbf{Req-T23-LC-07} & \kpimark{} & & \kpimark{} & & & \\
\textbf{Req-T23-LC-08} & & \kpimark{} & & & \kpimark{} & \\
\textbf{Req-T23-LC-09} & & \kpimark{} & & & \kpimark{} & \\
\textbf{Req-T23-LC-10} & & & & \kpimark{} & \kpimark{} & \\
\textbf{Req-T23-LC-11} & & \kpimark{} & \kpimark{} & & & \\
\textbf{Req-T23-LC-12} & & \kpimark{} & & \kpimark{} & & \\
\end{longtable}

\end{wptable}
